\answerAnchor{MATH1-CALC-0055}
\begin{solutionBox}
\textbf{A}\quad 交点由 \(1-x^2=0\) 得 \(x=\pm1\)。在区间内
\(1-x^2\ge0\)，故
\[
A=\int_{-1}^{1}(1-x^2)\mathrm dx
=\frac43.
\]
\end{solutionBox}
\begin{solutionBox}
\textbf{B}\quad 交点由 \(x^2=2-x\) 得 \(x=-2,1\)，直线在上：
\[
\begin{aligned}
A&=\int_{-2}^{1}(2-x-x^2)\mathrm dx\\
&=\left[2x-\frac{x^2}{2}-\frac{x^3}{3}\right]_{-2}^{1}
=\frac92.
\end{aligned}
\]
区间中点取值验证高度为正。
\end{solutionBox}
\begin{solutionBox}
\textbf{C}\quad \(r=1+\cos\theta\ge0\)，
\(0\le\theta\le2\pi\) 恰扫描心形线一次。故
\[
\begin{aligned}
A&=\frac12\int_0^{2\pi}(1+\cos\theta)^2\mathrm d\theta\\
&=\frac12(2\pi+\pi)=\frac{3\pi}{2}.
\end{aligned}
\]
\end{solutionBox}
\begin{solutionBox}
\textbf{M}\quad 两曲线相交时
\[
2a\cos\theta=a,\qquad \theta=\pm\frac\pi3.
\]
目标区域在 \(|\theta|\le\pi/3\) 内，外半径为 \(2a\cos\theta\)，
内半径为 \(a\)。于是
\[
\begin{aligned}
A&=\frac12\int_{-\pi/3}^{\pi/3}
\left(4a^2\cos^2\theta-a^2\right)\mathrm d\theta\\
&=a^2\left(\frac\pi3+\frac{\sqrt3}{2}\right).
\end{aligned}
\]
\end{solutionBox}

\answerAnchor{MATH1-CALC-0056}
\begin{solutionBox}
\textbf{A}\quad 外半径为 \(1\)，内半径为 \(x\)，故截面为垫片：
\[
V=\pi\int_0^1(1-x^2)\mathrm dx=\frac{2\pi}{3}.
\]
\end{solutionBox}
\begin{solutionBox}
\textbf{B}\quad 交点为 \(x=0,1\)。柱壳半径 \(x\)，高度
\(\sqrt x-x^2\)，故
\[
V=2\pi\int_0^1x(\sqrt x-x^2)\mathrm dx
=2\pi\left(\frac25-\frac14\right)
=\frac{3\pi}{10}.
\]
\end{solutionBox}
\begin{solutionBox}
\textbf{C}\quad 柱壳半径为 \(x+1\)，高度为 \(4-x^2\)，故
\[
V=2\pi\int_0^2(x+1)(4-x^2)\mathrm dx
=2\pi\cdot\frac{28}{3}
=\frac{56\pi}{3}.
\]
\end{solutionBox}
\begin{solutionBox}
\textbf{M}\quad 截面积为
\[
A(x)=x^2(1-x)^2.
\]
因此
\[
\begin{aligned}
V&=\int_0^1x^2(1-x)^2\mathrm dx\\
&=\int_0^1(x^2-2x^3+x^4)\mathrm dx
=\frac13-\frac12+\frac15=\frac1{30}.
\end{aligned}
\]
\end{solutionBox}

\answerAnchor{MATH1-CALC-0057}
\begin{solutionBox}
\textbf{A}\quad \(y'=2\)，故
\[
s=\int_0^3\sqrt{1+4}\mathrm dx=3\sqrt5.
\]
这也等于端点 \((0,0),(3,6)\) 的距离。
\end{solutionBox}
\begin{solutionBox}
\textbf{B}\quad
\[
y'=x\sqrt{x^2+2},
\]
\[
\sqrt{1+(y')^2}
=\sqrt{1+x^2(x^2+2)}
=x^2+1.
\]
\[
\begin{aligned}
s&=\int_0^1(x^2+1)\mathrm dx
=\frac43.
\end{aligned}
\]
\end{solutionBox}
\begin{solutionBox}
\textbf{C}\quad
\[
x'=2t,\qquad y'=t^2-1,
\]
\[
\sqrt{(x')^2+(y')^2}
=\sqrt{4t^2+(t^2-1)^2}
=t^2+1,
\]
\[
s=\int_{-1}^{1}(t^2+1)\mathrm dt=\frac83.
\]
\end{solutionBox}
\begin{solutionBox}
\textbf{M}\quad \(r'=e^\theta=r\)，故
\[
s=\int_0^{\ln2}\sqrt{r^2+(r')^2}\mathrm d\theta
=\sqrt2\int_0^{\ln2}e^\theta\mathrm d\theta
=\sqrt2.
\]
\end{solutionBox}

\answerAnchor{MATH1-CALC-0058}
\begin{solutionBox}
\textbf{A}\quad 半径恒为 \(2\)，弧长微元为 \(\mathrm dx\)，故
\[
S=2\pi\int_0^3 2\,\mathrm dx=12\pi.
\]
这与半径 \(2\)、高 \(3\) 的圆柱侧面积一致。
\end{solutionBox}
\begin{solutionBox}
\textbf{B}\quad
\[
y'=\frac1{2\sqrt x},\qquad
y\sqrt{1+(y')^2}=\frac12\sqrt{4x+1}.
\]
因此
\[
\begin{aligned}
S&=2\pi\int_0^1\frac12\sqrt{4x+1}\mathrm dx\\
&=\frac\pi6\left(5\sqrt5-1\right).
\end{aligned}
\]
零端虽有竖切线，但化简后的被积函数连续可积。
\end{solutionBox}
\begin{solutionBox}
\textbf{C}\quad 参数速度为 \(\sqrt{4t^2+1}\)，半径为 \(t\)，故
\[
S=2\pi\int_1^2t\sqrt{4t^2+1}\,\mathrm dt
=\frac\pi6\left(17\sqrt{17}-5\sqrt5\right).
\]
\end{solutionBox}
\begin{solutionBox}
\textbf{M}\quad 在给定区间上 \(\sin t,\cos t\ge0\)，且
\[
\sqrt{(x')^2+(y')^2}
=3\sin t\cos t.
\]
所以
\[
S=2\pi\int_0^{\pi/2}\sin^3t\cdot3\sin t\cos t\,\mathrm dt
=6\pi\int_0^{\pi/2}\sin^4t\cos t\,\mathrm dt
=\frac{6\pi}{5}.
\]
\end{solutionBox}
